% ========================================================================
% DIAGRAMY DO ROZDZIAŁU 4
% Wklej te diagramy w odpowiednie miejsca w chapter4_implementation.tex
% Wymagane pakiety w preambule:
%   \usepackage{tikz}
%   \usetikzlibrary{shapes.geometric, arrows.meta, positioning, fit,
%                    backgrounds, calc, decorations.pathreplacing}
%   \usepackage{float}
% ========================================================================


% ========================================================================
% DIAGRAM 1: Architektura trójwarstwowa
% Wstaw po sekcji 4.1 "Architektura rozwiązania"
% ========================================================================

\begin{figure}[H]
\centering
\begin{tikzpicture}[
    node distance=1.2cm and 2cm,
    box/.style={draw, rounded corners=4pt, minimum width=3.2cm, minimum height=1cm,
                font=\small, align=center, thick},
    layer/.style={draw, dashed, rounded corners=8pt, inner sep=12pt,
                  fill=#1, fill opacity=0.08},
    arrow/.style={-{Stealth[length=6pt]}, thick},
    doublearrow/.style={{Stealth[length=6pt]}-{Stealth[length=6pt]}, thick},
    label/.style={font=\footnotesize\bfseries, text=#1}
]

% === FRONTEND LAYER ===
\node[box, fill=blue!8] (react) {React 18.3\\SPA};
\node[box, fill=blue!8, right=1.5cm of react] (tailwind) {TailwindCSS\\UI Components};
\node[box, fill=blue!8, right=1.5cm of tailwind] (axios) {Axios\\HTTP Client};

\begin{scope}[on background layer]
    \node[layer=blue, fit=(react)(tailwind)(axios), label={[label=blue!70]above:Warstwa klienta (Frontend --- port 3000)}] (frontend) {};
\end{scope}

% === API LAYER ===
\node[box, fill=orange!12, below=2cm of tailwind] (rest) {REST API\\JSON / FormData};

% === BACKEND LAYER ===
\node[box, fill=green!8, below=2cm of rest, xshift=-3.5cm] (controllers) {Controllers\\Routing};
\node[box, fill=green!8, right=1cm of controllers] (services) {Services\\Logika biznesowa};
\node[box, fill=green!8, right=1cm of services] (optimizer) {ScheduleOptimizer\\SSGS / MORCPSP};
\node[box, fill=green!8, right=1cm of optimizer] (repos) {Repositories\\JPA / Hibernate};

\begin{scope}[on background layer]
    \node[layer=green, fit=(controllers)(services)(optimizer)(repos),
          label={[label=green!60!black]above:Warstwa serwera (Spring Boot 3.2.4 --- port 8080)}] (backend) {};
\end{scope}

% === DATABASE LAYER ===
\node[box, fill=red!8, below=2cm of services, xshift=0.5cm,
      minimum width=4cm] (db) {PostgreSQL\\Baza danych};
\node[box, fill=red!8, right=2cm of db] (hikari) {HikariCP\\Pula połączeń};

\begin{scope}[on background layer]
    \node[layer=red, fit=(db)(hikari),
          label={[label=red!60!black]above:Warstwa danych}] (database) {};
\end{scope}

% === ARROWS ===
\draw[doublearrow, blue!60] (axios) -- node[right, font=\scriptsize, text=gray] {HTTP/CORS\\JWT Cookies} (rest);
\draw[arrow, orange!60] (rest) -- (controllers);
\draw[arrow, green!60] (controllers) -- (services);
\draw[arrow, green!60] (services) -- (optimizer);
\draw[arrow, green!60] (services) -- (repos);
\draw[arrow, green!60] (repos) -- (hikari);
\draw[arrow, red!60] (hikari) -- (db);

% Side labels
\node[font=\tiny, text=gray, rotate=90, anchor=south] at ($(react.west)+(-0.3,0)$) {Port 3000};

\end{tikzpicture}
\caption{Architektura trójwarstwowa systemu FlowLink}
\label{fig:architecture}
\end{figure}


% ========================================================================
% DIAGRAM 2: Diagram encji (ERD uproszczony)
% Wstaw w sekcji 4.2.1 "Model danych"
% ========================================================================

\begin{figure}[H]
\centering
\begin{tikzpicture}[
    node distance=1.5cm and 3cm,
    entity/.style={draw, thick, rounded corners=3pt, minimum width=3cm,
                   minimum height=0.8cm, font=\small\bfseries, fill=#1},
    attr/.style={font=\tiny, text=gray!70!black, align=left},
    rel/.style={-{Stealth[length=5pt]}, thick, #1},
    manyrel/.style={{Diamond[length=5pt]}-{Stealth[length=5pt]}, thick, #1}
]

% Entities
\node[entity=blue!10] (user) {User};
\node[entity=green!10, right=4cm of user] (project) {Project};
\node[entity=orange!10, below=3cm of project] (task) {Task};
\node[entity=purple!10, below=2.5cm of user] (comment) {Comment};
\node[entity=red!10, left=1cm of comment, yshift=-1.5cm] (notification) {Notification};
\node[entity=yellow!15, right=4cm of task] (refreshtoken) {RefreshToken};

% Attributes
\node[attr, below=0.15cm of user] {id, email, firstname\\lastname, password};
\node[attr, below=0.15cm of project] {id, projectKey, summary\\nextTaskNumber};
\node[attr, below=0.15cm of task] {id, taskNumber, summary\\status, startDate, dueDate\\priority, progress, assignee\_id};
\node[attr, below=0.15cm of comment] {id, content, author\_id\\task\_id, parentComment\_id};

% Relationships
\draw[rel=blue!60] (user) -- node[above, font=\tiny] {owner (1:N)} (project);
\draw[manyrel=blue!40, dashed] (user) -- node[below, font=\tiny, sloped] {members (M:N)} (project);
\draw[rel=green!60] (project) -- node[right, font=\tiny] {tasks (1:N)} (task);
\draw[rel=blue!60, bend right=20] (user.south) to node[left, font=\tiny] {assignee} (task.west);
\draw[rel=orange!60] (task) -- node[below, font=\tiny, sloped] {comments (1:N)} (comment);
\draw[rel=blue!60] (user) -- node[left, font=\tiny] {author} (comment);
\draw[rel=blue!60, bend right=25] (user.south west) to node[left, font=\tiny] {recipient (1:N)} (notification);

% Self-reference: task dependencies
\draw[manyrel=orange!60, looseness=4] (task.south east) to[out=-30, in=30]
    node[right, font=\tiny] {dependencies\\(M:N)} (task.north east);

% Self-reference: comment threading
\draw[rel=purple!40, looseness=5, dashed] (comment.south) to[out=-150, in=150]
    node[left, font=\tiny] {parentComment} (comment.north);

\end{tikzpicture}
\caption{Uproszczony diagram związków encji (ERD) systemu FlowLink}
\label{fig:er_diagram}
\end{figure}


% ========================================================================
% DIAGRAM 3: Schemat blokowy algorytmu SSGS
% Wstaw w sekcji 4.3.1 po opisie faz algorytmu
% ========================================================================

\begin{figure}[H]
\centering
\begin{tikzpicture}[
    node distance=1cm,
    startstop/.style={draw, rounded corners=8pt, minimum width=3.5cm, minimum height=0.8cm,
                      fill=slate!15, font=\small, thick, align=center},
    process/.style={draw, minimum width=4.5cm, minimum height=0.8cm,
                    fill=blue!6, font=\small, thick, align=center},
    decision/.style={draw, diamond, aspect=2.5, minimum width=2cm,
                     fill=orange!8, font=\small, thick, align=center, inner sep=1pt},
    io/.style={draw, trapezium, trapezium left angle=70, trapezium right angle=110,
               minimum width=3.5cm, fill=green!6, font=\small, thick, align=center},
    arrow/.style={-{Stealth[length=5pt]}, thick},
    phase/.style={font=\scriptsize\bfseries, text=gray}
]

% Start
\node[startstop] (start) {Wejście: Lista TaskDTO,\\horyzont $T_0$, $\alpha$, $\beta$};

% Phase 1
\node[process, below=0.8cm of start] (build) {Budowa grafu zadań\\(predecessors, successors)};
\node[phase, left=0.3cm of build, anchor=east] {Faza 1};

% Phase 2
\node[process, below=0.8cm of build] (topo) {Oblicz $|\text{Succ}_j^*|$\\(odwrotny porządek topologiczny)};
\node[process, below=0.6cm of topo] (sort) {Sortuj wg $score_j = w_j \cdot \frac{H-d_j}{H} + 0{,}5 \cdot |\text{Succ}_j^*|$};
\node[phase, left=0.3cm of topo, anchor=east] {Faza 2};

% Phase 3 - loop
\node[io, below=0.8cm of sort] (next) {Weź następne zadanie $j$\\z listy priorytetowej};
\node[phase, left=0.3cm of next, anchor=east] {Faza 3};

\node[process, below=0.8cm of next] (prec) {$ES_j = \max_{i \in Pred_j} C_i$};

\node[decision, below=0.8cm of prec] (hasres) {Assignee\\$\neq$ null?};

\node[process, right=2.5cm of hasres] (resource) {Szukaj wolnego slotu\\w BitSet[assignee]};

\node[process, below=1.2cm of hasres] (place) {$S_j = \max(0, ES_j, \text{slot})$\\$C_j = S_j + p_j$\\Oznacz BitSet $[S_j, C_j)$};

\node[decision, below=0.8cm of place] (more) {Pozostały\\zadania?};

% Phase 4
\node[process, below=1cm of more] (metrics) {Oblicz: $\sum w_j T_j$, $C_{max}$,\\$Z = \alpha \sum w_j T_j + \beta C_{max}$};
\node[phase, left=0.3cm of metrics, anchor=east] {Faza 4};

\node[startstop, below=0.8cm of metrics] (end) {Wynik: ScheduleResult\\(sugestie, metryki)};

% Arrows
\draw[arrow] (start) -- (build);
\draw[arrow] (build) -- (topo);
\draw[arrow] (topo) -- (sort);
\draw[arrow] (sort) -- (next);
\draw[arrow] (next) -- (prec);
\draw[arrow] (prec) -- (hasres);
\draw[arrow] (hasres) -- node[above, font=\scriptsize] {Tak} (resource);
\draw[arrow] (hasres) -- node[left, font=\scriptsize] {Nie} (place);
\draw[arrow] (resource) |- (place);
\draw[arrow] (place) -- (more);
\draw[arrow] (more) -- node[left, font=\scriptsize] {Nie} (metrics);
\draw[arrow] (more.east) -- ++(2,0) |- node[near start, above, font=\scriptsize] {Tak} (next.east);
\draw[arrow] (metrics) -- (end);

\end{tikzpicture}
\caption{Schemat blokowy algorytmu SSGS zaimplementowanego w~klasie \texttt{ScheduleOptimizer}}
\label{fig:ssgs_flowchart}
\end{figure}


% ========================================================================
% DIAGRAM 4: Przepływ uwierzytelniania JWT
% Wstaw w sekcji 4.2.2 "Mechanizm uwierzytelniania"
% ========================================================================

\begin{figure}[H]
\centering
\begin{tikzpicture}[
    node distance=0.6cm,
    actor/.style={draw, thick, rounded corners=3pt, minimum width=2.2cm,
                  minimum height=0.7cm, font=\small\bfseries, fill=#1},
    msg/.style={-{Stealth[length=5pt]}, thick, #1},
    msglabel/.style={font=\scriptsize, fill=white, inner sep=1pt},
    timeline/.style={thick, gray!40},
    note/.style={draw, dashed, rounded corners=3pt, fill=yellow!5,
                 font=\scriptsize, align=left, text width=3.5cm}
]

% Actors
\node[actor=blue!10] (client) {Klient (React)};
\node[actor=green!10, right=5cm of client] (server) {Serwer (Spring)};
\node[actor=red!10, right=4cm of server] (db) {PostgreSQL};

% Timelines
\draw[timeline] (client.south) -- ++(0,-11.5);
\draw[timeline] (server.south) -- ++(0,-11.5);
\draw[timeline] (db.south) -- ++(0,-11.5);

% Step 1: Login
\draw[msg=blue!60] ($(client.south)+(0,-0.8)$) -- node[msglabel, above] {1. POST /api/auth/login} ($(server.south)+(0,-0.8)$);

% Step 2: Verify
\draw[msg=green!60] ($(server.south)+(0,-1.5)$) -- node[msglabel, above] {2. SELECT user WHERE email=?} ($(db.south)+(0,-1.5)$);
\draw[msg=red!40, dashed] ($(db.south)+(0,-2.0)$) -- node[msglabel, above] {User entity} ($(server.south)+(0,-2.0)$);

% Step 3: BCrypt
\node[note, anchor=west] at ($(server.south)+(0.3,-2.8)$) {3. BCrypt.matches(\\password, hash)\\Generuj JWT access (15min)\\+ refresh token (7 dni)};

% Step 4: Response with cookies
\draw[msg=green!60] ($(server.south)+(0,-4.0)$) -- node[msglabel, above] {4. Set-Cookie: accessToken, refreshToken} ($(client.south)+(0,-4.0)$);
\node[msglabel, below] at ($(client.south)!0.5!(server.south)+(0,-4.0)$) {\textit{HttpOnly, Secure, SameSite=Lax}};

% Step 5: Authenticated request
\draw[msg=blue!60] ($(client.south)+(0,-5.5)$) -- node[msglabel, above] {5. GET /api/projects (Cookie: accessToken)} ($(server.south)+(0,-5.5)$);

% Step 6: JwtFilter validates
\node[note, anchor=west] at ($(server.south)+(0.3,-6.3)$) {6. JwtAuthenticationFilter:\\validateToken(accessToken)\\userId $\rightarrow$ request attribute};

% Step 7: 401 + Refresh
\draw[msg=orange!60, dashed] ($(server.south)+(0,-7.5)$) -- node[msglabel, above] {7. 401 Unauthorized (token expired)} ($(client.south)+(0,-7.5)$);

\draw[msg=blue!60] ($(client.south)+(0,-8.3)$) -- node[msglabel, above] {8. POST /api/auth/refresh (Cookie: refreshToken)} ($(server.south)+(0,-8.3)$);

% Step 9: New tokens
\draw[msg=green!60] ($(server.south)+(0,-9.0)$) -- node[msglabel, above] {9. Set-Cookie: nowy accessToken} ($(client.south)+(0,-9.0)$);

% Step 10: Retry
\draw[msg=blue!60] ($(client.south)+(0,-9.8)$) -- node[msglabel, above] {10. Ponowienie oryginalnego żądania} ($(server.south)+(0,-9.8)$);
\draw[msg=green!60] ($(server.south)+(0,-10.5)$) -- node[msglabel, above] {11. 200 OK + dane} ($(client.south)+(0,-10.5)$);

% Axios note
\node[note, anchor=east, text width=2.8cm] at ($(client.south)+(-0.3,-8.5)$) {Axios interceptor:\\kolejkuje żądania\\podczas odświeżania\\(1 refresh naraz)};

\end{tikzpicture}
\caption{Diagram sekwencji przepływu uwierzytelniania JWT z~automatycznym odświeżaniem tokena}
\label{fig:jwt_sequence}
\end{figure}


% ========================================================================
% DIAGRAM 5: Komponenty modułu optymalizacyjnego
% Wstaw na początku sekcji 4.3 "Moduł optymalizacyjny"
% ========================================================================

\begin{figure}[H]
\centering
\begin{tikzpicture}[
    node distance=1cm and 1.5cm,
    comp/.style={draw, thick, rounded corners=4pt, minimum width=3.5cm,
                 minimum height=1.2cm, font=\small, align=center, fill=#1},
    api/.style={draw, thick, minimum width=3cm, minimum height=0.6cm,
                font=\scriptsize\ttfamily, fill=orange!8, rounded corners=2pt},
    arrow/.style={-{Stealth[length=5pt]}, thick},
    dashedarrow/.style={-{Stealth[length=5pt]}, thick, dashed, gray},
    group/.style={draw, dashed, rounded corners=8pt, inner sep=10pt, fill=#1, fill opacity=0.05}
]

% Backend components
\node[comp=green!10] (controller) {OptimizationController\\{\scriptsize @RestController}};
\node[comp=green!10, below=1.2cm of controller] (service) {OptimizationService\\{\scriptsize @Service, @Transactional}};
\node[comp=blue!10, below left=1.2cm and 0.5cm of service] (optimizer) {ScheduleOptimizer\\{\scriptsize Stateless, SSGS}};
\node[comp=blue!10, below right=1.2cm and 0.5cm of service] (cpm) {CriticalPathMethodHelper\\{\scriptsize Forward/Backward Pass}};
\node[comp=gray!10, right=1.5cm of service] (taskrepo) {TaskRepository\\{\scriptsize Entity Graph: withDetails}};
\node[comp=gray!10, above=0.5cm of taskrepo] (projectrepo) {ProjectRepository\\{\scriptsize findByProjectKey}};

% API endpoints
\node[api, above left=0.8cm and -0.5cm of controller] (simulate) {POST /simulate};
\node[api, above right=0.8cm and -0.5cm of controller] (apply) {POST /apply};

% Frontend components
\node[comp=purple!10, above=2.5cm of controller, xshift=-3cm] (projpage) {ProjectsPage\\{\scriptsize handleOptimize()}};
\node[comp=purple!10, above=2.5cm of controller, xshift=3cm] (metrics) {OptimizationMetrics\\{\scriptsize Panel wyników}};
\node[comp=purple!10, below=0.3cm of metrics] (timeline) {TimelineView\\{\scriptsize Ghost Bars}};

% Groups
\begin{scope}[on background layer]
    \node[group=green, fit=(controller)(service)(optimizer)(cpm)(taskrepo)(projectrepo),
          label={[font=\footnotesize\bfseries, text=green!50!black]below:Backend (Java / Spring Boot)}] {};
    \node[group=purple, fit=(projpage)(metrics)(timeline),
          label={[font=\footnotesize\bfseries, text=purple!50!black]above:Frontend (React)}] {};
\end{scope}

% Arrows
\draw[arrow] (simulate) -- (controller);
\draw[arrow] (apply) -- (controller);
\draw[arrow] (controller) -- (service);
\draw[arrow] (service) -- (optimizer);
\draw[arrow] (service) -- (cpm);
\draw[arrow] (service) -- (taskrepo);
\draw[dashedarrow] (service) -- (projectrepo);

\draw[arrow, purple!60] (projpage) -- node[left, font=\scriptsize] {API call} (simulate);
\draw[arrow, purple!60, dashed] (projpage) -- (metrics);
\draw[arrow, purple!60, dashed] (projpage) -- (timeline);

% DTOs
\node[font=\tiny, text=gray, rotate=-90, anchor=north] at ($(controller)!0.5!(service)+(-1.8,0)$) {OptimizationResultDTO};
\node[font=\tiny, text=gray] at ($(service)!0.5!(optimizer)+(0,0.3)$) {List<TaskDTO>};
\node[font=\tiny, text=gray] at ($(service)!0.5!(cpm)+(0,0.3)$) {isCritical};

\end{tikzpicture}
\caption{Diagram komponentów modułu optymalizacyjnego i~jego integracja z~interfejsem użytkownika}
\label{fig:optimization_components}
\end{figure}


% ========================================================================
% DIAGRAM 6: Przepływ optymalizacji z perspektywy użytkownika
% Wstaw w sekcji 4.4.2 po opisie Ghost Bars
% ========================================================================

\begin{figure}[H]
\centering
\begin{tikzpicture}[
    node distance=0.8cm,
    step/.style={draw, thick, rounded corners=4pt, minimum width=5cm,
                 minimum height=0.9cm, font=\small, align=center, fill=#1},
    arrow/.style={-{Stealth[length=5pt]}, thick},
    choice/.style={draw, diamond, aspect=2, font=\small, thick,
                   fill=orange!8, inner sep=2pt, align=center}
]

\node[step=blue!8] (s1) {1. Użytkownik klika\\``Optimize'' na osi czasu};
\node[step=gray!8, below=0.7cm of s1] (s2) {2. Frontend wysyła\\POST /api/optimization/simulate};
\node[step=green!8, below=0.7cm of s2] (s3) {3. Backend: pobranie zadań,\\CPM, SSGS optymalizacja};
\node[step=gray!8, below=0.7cm of s3] (s4) {4. Odpowiedź: sugestie +\\metryki ``przed'' i ``po''};
\node[step=purple!8, below=0.7cm of s4] (s5) {5. Wyświetlenie Ghost Bars\\+ panel metryk};
\node[choice, below=0.8cm of s5] (s6) {Użytkownik\\decyduje};
\node[step=green!12, below left=1cm and 1cm of s6] (s7a) {``Apply Schedule''\\POST /apply → zapis dat};
\node[step=red!8, below right=1cm and 1cm of s6] (s7b) {``Dismiss''\\usunięcie duchów};
\node[step=blue!8, below=0.7cm of s7a] (s8) {Odświeżenie widoku\\z nowymi datami};

\draw[arrow] (s1) -- (s2);
\draw[arrow] (s2) -- (s3);
\draw[arrow] (s3) -- (s4);
\draw[arrow] (s4) -- (s5);
\draw[arrow] (s5) -- (s6);
\draw[arrow] (s6) -| node[above left, font=\scriptsize] {Akceptuj} (s7a);
\draw[arrow] (s6) -| node[above right, font=\scriptsize] {Odrzuć} (s7b);
\draw[arrow] (s7a) -- (s8);

\end{tikzpicture}
\caption{Przepływ interakcji użytkownika z~modułem optymalizacji harmonogramu}
\label{fig:optimization_flow}
\end{figure}
